\section{The CIDG Environment}
\label{sec:env}

% [F2 content: environment + simulator. Source: sim/engine.py, ARCHITECTURE.md]
\paragraph{Scenario specs as a single source of truth.}
Each incident is a CIDG scenario: a typed dependency graph with a root-cause fault, a
propagation rule, a victim service on which a uniform alert fires, and a labelled trap
remediation. Specs are authored in a small DSL (\texttt{sim/spec.py},
\texttt{scenarios/cidg/*.yaml}) and seeded from real public post-mortems
(\texttt{opensre-traj/specs/real/*.json}: GCP Service Control, Kinesis, Cloudflare,
Slack/Consul, GitHub, Datadog). We release 51 generated scenarios.

\paragraph{Tier-A: deterministic in-process simulator.}
\texttt{sim/engine.py} runs \texttt{propagate()} over the dependency graph so that the cascade
is \emph{emergent}, not scripted. A root-cause-aware \texttt{is\_resolved} distinguishes a
true fix from a symptom suppression, and trap detection flags the locally-obvious-but-harmful
action. This tier is free and reproducible.

\paragraph{Tier-B: M-real live cluster.}
For high fidelity, \texttt{mreal/server.py} on GKE implements a real-RPS call mesh where
services call upstreams over HTTP, so cascades are \emph{physically} emergent. A
kube-prometheus-stack (Prometheus, Alertmanager, Grafana) fires the same uniform alert on the
victim, mirroring production observability.

\paragraph{Rollout harness.}
\texttt{rex/harness.py} and \texttt{rex/loop.py} drive the agent loop:
\(\text{observe} \rightarrow \text{propose tool} \rightarrow \text{apply} \rightarrow
\text{re-diagnose}\), emitting FIREBALL-schema trajectories
(\(\text{state\_before} \rightarrow \text{tool} \rightarrow \text{state\_after} \rightarrow
\text{reward}\)).
